\documentclass[../cours_tle.tex]{subfiles}
\begin{document}
\setcounter{chapter}{7}
\chapter{Sphères, barycentres et mesures}
\label{ch:8}
\begin{objectif}
Spécialité pour le plan médiateur et les mesures ; approfondissements pour les sphères et les barycentres. Consolider les lieux géométriques et préparer l'algèbre affine. Prérequis : chapitres 5 à 7.
\end{objectif}
\section{Sphère et plan médiateur}
Une sphère de centre $C$ et de rayon $R>0$ est l'ensemble des points $M$ tels que $CM=R$. En repère orthonormé :
$(x-x_C)^2+(y-y_C)^2+(z-z_C)^2=R^2$.
Pour trouver son intersection avec une droite, substituer sa représentation paramétrique : une équation du second degré donne zéro, une ou deux intersections réelles.

Les points équidistants de $A\ne B$ satisfont $MA^2=MB^2$. En développant autour du milieu $I$,
$MA^2-MB^2=2\overrightarrow{IM}\cdot\overrightarrow{AB}$.
Le lieu est donc le plan passant par $I$ et normal à $\overrightarrow{AB}$.
\section{Intersection d'une sphère et d'un plan}
Soit $H$ le projeté de $C$ sur le plan. Pour $M$ dans ce plan,
$CM^2=CH^2+HM^2$. Si $d=CH$, l'intersection est vide pour $d>R$, réduite à $H$ pour $d=R$, un cercle de centre $H$ et rayon $\sqrt{R^2-d^2}$ pour $d<R$.
Le plan tangent en $T$ sur la sphère est normal à $\overrightarrow{CT}$ : Pythagore montre qu'il n'a que $T$ en commun avec la sphère.
Pour un tétraèdre non aplati, chercher un centre équidistant des quatre sommets conduit à trois équations linéaires obtenues par soustraction des carrés. Leurs normales indépendantes assurent l'unicité ; le rayon est ensuite la distance à un sommet.
\section{Barycentres et associativité}
\begin{definition}
Pour des points $A_i$ et des poids réels $\alpha_i$ de somme $s\ne0$, le barycentre $G$ est défini par $\sum_i\alpha_i\overrightarrow{GA_i}=0$.
\end{definition}
\begin{demonstration}[Existence, unicité et coordonnées]
Chasles transforme cette relation en
$\sum_i\alpha_i\overrightarrow{OA_i}-s\overrightarrow{OG}=0$.
Ainsi $\overrightarrow{OG}=s^{-1}\sum_i\alpha_i\overrightarrow{OA_i}$ détermine un unique point, indépendamment de l'origine.
\end{demonstration}
Regrouper des points dont la somme des poids est non nulle, puis les remplacer par leur barycentre affecté de cette somme, ne change pas la somme vectorielle : c'est l'associativité.
Le milieu correspond aux poids $(1,1)$ ; le centre de gravité du triangle aux poids $(1,1,1)$. Des poids négatifs peuvent placer le barycentre hors de l'enveloppe des points.
\section{Deux identités de Leibniz}
Pour tout $M$, Chasles et la définition donnent
$\sum_i\alpha_i\overrightarrow{MA_i}=s\overrightarrow{MG}$.
En développant $\|\overrightarrow{MG}+\overrightarrow{GA_i}\|^2$ et en sommant, les termes croisés s'annulent :
\[
\sum_i\alpha_i MA_i^2=s\,MG^2+\sum_i\alpha_i GA_i^2.
\]
Si les poids sont positifs, cette somme est minimale uniquement en $G$.
C'est déjà une méthode des moindres carrés.
\section{Aires et volumes : ne pas oublier la hauteur}
Pour un triangle engendré par $\vec u,\vec v$, sa hauteur relative à $\vec u$ est la norme de la composante de $\vec v$ orthogonale à $\vec u$. Le chapitre 6 donne
\[
\mathcal A=\frac12\sqrt{\|\vec u\|^2\|\vec v\|^2-(\vec u\cdot\vec v)^2}.
\]
Le volume du tétraèdre est $\mathcal A h/3$, avec $h$ distance du quatrième sommet au plan de base. Pour expliquer le facteur $1/3$, les sections parallèles à la base à hauteur $t$ depuis le sommet ont aire $\mathcal A(t/h)^2$ par similitude. Leur intégrale de $0$ à $h$ vaut $\mathcal A h/3$ (chapitre 19 ; interprétation du volume par sections admise).
\begin{finchapitre}
\section{Utilité et suite du parcours}
Barycentres et identités quadratiques reviendront en statistique, en optimisation et en géométrie affine. Les sections de solides préparent les intégrales multiples.
\begin{retenir}
Une longueur se mesure perpendiculairement à la base. Toujours vérifier que la somme des poids d'un barycentre est non nulle.
\end{retenir}
\end{finchapitre}
\end{document}
